\chapter{AI Agent Usage Guide}
\ChapterLabel{AIAgentGuide}
\index{AI agent}
\index{AI agent!context routing}
\index{AI agent!example prompts}
\index{AI agent!human review}
\index{AI agent!TikZ requirement}
\index{Appendix!AI Agent Usage Guide}

\section{How to Use This Appendix}
This appendix is the entry point for an AI agent that creates, reviews, or
maintains this architecture specification. The agent shall read this file
first, identify the requested task in
\hyperref[Table:AITaskRouting]{the task-routing table}, and load only the
authoritative sources listed for that task. It shall expand the context when a
referenced dependency, cross-cutting decision, or inconsistency requires it.

This appendix replaces substantive AI-oriented Markdown instructions. A thin
repository bootstrap may point an AI tool to this appendix, but its routing
tables, context packs, workflows, prompts, and checklists are maintained only
as LaTeX so they remain visible in the architecture book and versioned with the
specification.

\section{Governing Principles}
\index{source of truth}
The authored LaTeX files and editable diagram sources are the source of truth.
TikZ is the required source format for diagrams.
The generated PDF is a review artifact and shall not be edited directly. The
agent shall preserve one authoritative location for each fact and use links
instead of duplicating content across chapters.

An AI agent is a drafting, analysis, consistency-checking, and build assistant.
It shall use supplied project evidence, identify missing information, expose
material assumptions, and distinguish proposals from approved decisions. It
shall not invent requirements approval, implementation status, test results,
verification evidence, validation evidence, waivers, or release authorization.

\section{AI Quick Start}
\label{Section:AIQuickStart}
\index{AI agent!quick start}

\begin{enumerate}
\item Read this appendix and identify the requested outcome.
\item Before creating or renaming a project-owned file, ask for or confirm the
project file namespace. Do not invent it; use the confirmed namespace for every
file created during the task.
\item Inspect \hyperref[Table:DomainProfileApplicability]{the domain-profile
applicability table} and \hyperref[Table:ProjectDomainProfiles]{the project
profile decisions}. Ask an authorized human to resolve any unreviewed profile
selection. Load every Applicable or TBD profile affected by the task.
\item Select the closest entry in
\hyperref[Table:AITaskRouting]{the task-routing table}.
\item Read the listed primary files and the applicable context pack. Read
\texttt{tgmManual.tex} when document metadata, macros, inclusion order, or
build structure is affected.
\item Inspect available project evidence and report conflicts, missing facts,
and consequential assumptions. Ask for a human decision when an answer would
materially change scope, architecture, safety, security, cost, schedule, or
release readiness.
\item Make narrowly scoped changes in the authoritative files. Preserve stable
IDs and update every affected reciprocal mapping.
\item Add a linked document-history entry for an approved document change and
record material AI interactions in \ChapterRef{AIPromptLog}.
\item Run the complete build, inspect the log, render affected PDF pages, and
complete \hyperref[Table:AICompletionChecklist]{the completion checklist}.
\end{enumerate}

\section{Repository AI Bootstrap}
\label{Section:AIRepositoryBootstrap}
\index{AI agent!repository bootstrap}
\index{AGENTS.md}
\index{CLAUDE.md}
\index{SKILL.md}

The LaTeX architecture is authoritative, but an AI tool may require a
repository-level discovery file before it knows where to find these
instructions. Create a thin bootstrap file for each AI tool that the project
actually uses. The bootstrap shall point to this chapter and the relevant
LaTeX sources; it shall not reproduce their requirements, workflows, or
architectural decisions. When policy changes, update the authoritative LaTeX
source rather than copying the change into each bootstrap file.

Tool-mandated discovery names are documented exceptions to the project file
namespace requirement. This exception applies only to the exact filename or
location required for discovery. It does not exempt files created while
performing the instructed work from the namespace requirements in
\hyperref[Table:DocumentAuthoringRequirements]{the document-authoring
requirements}.

Create only the applicable files:

\footnotesize
\begin{tabularx}{\linewidth}{L{4.5cm}Y}
\caption{Repository AI Bootstrap Files\label{Table:AIRepositoryBootstrapFiles}}\\
\toprule
\textbf{File} & \textbf{When to Create It}\\
\midrule
\endfirsthead
\multicolumn{2}{l}{\small\itshape Repository AI Bootstrap Files (continued)}\\
\toprule
\textbf{File} & \textbf{When to Create It}\\
\midrule
\endhead
\bottomrule
\endlastfoot
\texttt{AGENTS.md} & The repository is maintained with Codex or another agent
that discovers this filename.\\
\texttt{CLAUDE.md} & The repository is maintained with Claude Code.\\
\texttt{.github/copilot-}\newline\texttt{instructions.md} & The repository uses repository-wide
GitHub Copilot instructions.\\
\texttt{SKILL.md} & The repository is intentionally packaged as, or contains,
a reusable AI skill. Do not create this file merely as a second copy of
\texttt{AGENTS.md}.\\
\end{tabularx}
\normalsize

Each bootstrap file shall identify the root architecture source, point to this
chapter, require task-appropriate context loading through
\hyperref[Table:AITaskRouting]{the task-routing table}, direct architecture
policy changes back to LaTeX, and point to \ChapterRef{DocumentGuide} for
building and verification. It shall also state the confirmed project file
namespace or direct the agent to obtain it before creating project-owned files.
Use the following content as a minimal tool-neutral pattern, adapting only the
discovery-specific wrapper:

\begin{verbatim}
# AI Instructions

The authoritative software architecture is maintained in
tgmManual.tex. Before architecture-related work, read
tgmAIAgentGuide.tex and use its task-routing table to identify the
additional LaTeX chapters required for the task.

Treat the authored LaTeX and editable diagram sources as the source
of truth. Do not duplicate their requirements, conventions,
workflows, or architectural decisions in this file. When policy
changes, update the authoritative LaTeX source.

Before creating project-owned files, obtain or confirm the project
namespace required by tgmDocumentGuide.tex. Follow that guide's
build and verification instructions before completing the task.
\end{verbatim}

\section{Task Routing}
\index{AI agent!task routing}
The file lists below are minimum starting contexts, not permission to ignore a
dependency discovered during the task. A task that changes an artifact shall
also inspect every register that contains a reciprocal mapping to that
artifact.

\small
\begin{tabularx}{\linewidth}{L{4.0cm}Y}
\caption{AI Task-Routing Table\label{Table:AITaskRouting}}\\
\toprule
\textbf{Task} & \textbf{Minimum Authoritative Sources}\\
\midrule
\endfirsthead
\multicolumn{2}{l}{\small\itshape AI Task-Routing Table (continued)}\\
\toprule
\textbf{Task} & \textbf{Minimum Authoritative Sources}\\
\midrule
\endhead
\midrule
\multicolumn{2}{r}{\small\itshape Continued on next page}\\
\endfoot
\bottomrule
\endlastfoot
Initialize a project & \texttt{tgmManual.tex}, \texttt{prologue.tex}, \texttt{tgmIntroduction.tex}, \texttt{tgmDevelopmentPlan.tex}, \texttt{tgmSystemContext.tex}, each proposed or selected domain profile, and the supplied project evidence.\\
Tailor an embedded system & \texttt{tgmEmbeddedSystemsProfile.tex},
\texttt{tgmDevelopmentPlan.tex}, \texttt{tgmSystemContext.tex},
\texttt{tgmRequirements.tex}, the affected architecture and interface
chapters, \texttt{tgmTests.tex}, \texttt{tgmTraceability.tex}, and the
controlled hardware and project evidence.\\
Create repository AI bootstrap files & \texttt{tgmAIAgentGuide.tex},
\texttt{tgmDocumentGuide.tex}, \texttt{tgmManual.tex}, evidence identifying the
AI tools used by the project, and each tool's required discovery filename.\\
Develop user stories & \texttt{tgmSystemContext.tex}, \texttt{tgmUserStories.tex}, \texttt{tgmRequirements.tex}, and stakeholder evidence.\\
Add or change requirements & \texttt{tgmUserStories.tex}, \texttt{tgmRequirements.tex}, every Applicable or TBD domain profile, \texttt{tgmUseCases.tex}, \texttt{tgmTests.tex}, \texttt{tgmTraceability.tex}, \texttt{tgmVerification.tex}, and \texttt{tgmGlossary.tex}.\\
Add or change use cases & \texttt{tgmUserStories.tex}, \texttt{tgmRequirements.tex}, \texttt{tgmUseCases.tex}, \texttt{tgmTests.tex}, and \texttt{tgmTraceability.tex}.\\
Develop or change architecture & \texttt{tgmSystemContext.tex}, \texttt{tgmRequirements.tex}, \texttt{tgmArchitectureOverview.tex}, the affected interface, logical, process, development, physical, or implementation chapters, and their \texttt{TikZ/} diagram sources.\\
Add or change tests & \texttt{tgmRequirements.tex}, \texttt{tgmUseCases.tex}, \texttt{tgmTests.tex}, \texttt{tgmTraceability.tex}, \texttt{tgmVerification.tex}, and \texttt{tgmTestEvidence.tex}.\\
Perform verification & \texttt{tgmRequirements.tex}, \texttt{tgmTests.tex}, \texttt{tgmTraceability.tex}, \texttt{tgmVerification.tex}, and controlled test evidence.\\
Perform validation & \texttt{tgmUserStories.tex}, \texttt{tgmUseCases.tex}, \texttt{tgmValidation.tex}, and controlled stakeholder evidence.\\
Prepare a release & \texttt{tgmRequirements.tex}, \texttt{tgmTraceability.tex}, \texttt{tgmVerification.tex}, \texttt{tgmValidation.tex}, \texttt{tgmRelease.tex}, and controlled release evidence.\\
Review consistency & All registers containing the affected artifact types, \texttt{tgmGlossary.tex}, \texttt{tgmDocumentHistory.tex}, and \texttt{tgmManual.tex}.\\
Compile or troubleshoot & \texttt{tgmManual.tex}, \texttt{tgmDocumentGuide.tex}, the affected chapter, its diagram sources, and the build log.\\
\end{tabularx}
\normalsize

\section{Task-Specific Context Packs}
\index{AI agent!context pack}
A context pack is a reading boundary within this appendix, not a separate file.
Use the applicable pack together with the task-routing entry.

\subsection{Domain-Profile Context}
Determine the Applicable, Not Applicable, or TBD status of every available
domain profile during initialization. An authorized human shall approve the
selection and rationale. For later tasks, load every Applicable or TBD profile
whose concerns intersect the requested change. A profile adds review concerns
and artifact expectations; project-specific facts remain in their authoritative
core chapters. When the Embedded Systems profile applies, follow
\ChapterRef{EmbeddedSystemsProfile} and add hardware, firmware, timing,
resource, power, physical-interface, safety, security, manufacturing, update,
and field-service concerns to the task context. For requirements work, apply
the controlled translation rule in
Section~\ref{Section:DomainProfileRequirements}: each applicable profile
concern shall map to a requirement, an existing requirement, or an approved
Not Applicable decision with rationale.

\subsection{Requirements Context}
Establish the system boundary and stakeholder need before proposing a
requirement. Keep each requirement singular, feasible, unambiguous, and
verifiable. Apply the identifier grammar and family definitions in
Section~\ref{Section:RequirementIdentifierConvention}. Classify its primary
family or subtype, write binding behavior with
\textbf{shall}, assign one \gls{must}, \gls{should}, \gls{could}, or
\gls{would} priority, assign one lifecycle status, and map it to at least one
use case and one test case. A requirement may be marked Done only when its
mapped tests pass, verification is complete, traceability is current, and a
controlled implementation version is known. Evaluate every Applicable or TBD
domain profile that intersects the proposed change and follow
Section~\ref{Section:DomainProfileRequirements} before completing the
requirements work.

\subsection{Use-Case Context}
Start from an actor goal and a system boundary. Give each use case its own
table, identify participating actors, preconditions, trigger, success outcome,
main flow, alternate and exception flows, postconditions, mapped requirements,
and mapped tests. Describe externally observable interaction rather than
internal implementation detail. Maintain reciprocal mappings to every related
requirement.

\subsection{Architecture Context}
Trace each architectural decision to requirements, quality attributes,
constraints, risks, and interfaces. Identify alternatives and explain the
trade-off that selected the proposed structure. Keep the context, logical,
process, development, physical, interface, and implementation views consistent.
Use TikZ for diagrams.
Store each TikZ source under its owning chapter as required by
Section~\ref{Section:TikZAuthoring}, include it from that chapter, and compile
the PDF.

\subsection{Testing and Evidence Context}
Derive tests from requirements and use cases. Define preconditions, inputs,
procedure, expected result, and retained evidence. A test definition is not a
test result. Verification records objective compliance with requirements;
validation records whether stakeholder needs are satisfied in intended use.
Never create passing evidence or approval from an expected result.

\subsection{Release Context}
Review open, in-progress, done, and retired requirements; unresolved defects;
verification and validation results; operational readiness; migration and
rollback provisions; known limitations; and required approvals. Report a gap
instead of treating missing evidence as acceptance. Only an authorized human
may approve release.

\subsection{Document-Authoring Context}
Follow \ChapterRef{DocumentGuide} for file naming, LaTeX conventions, diagrams,
installation, and builds. Update the glossary and index when introducing a
controlled term. Update document history with a link to the changed location.
Before creating or renaming a project-owned file, ask for or confirm the
project file namespace and apply it consistently in accordance with
\hyperref[Table:DocumentAuthoringRequirements]{the document-authoring
requirements}.
Every TikZ source shall satisfy the path, inclusion, caption, label, and
reference requirements in Section~\ref{Section:TikZAuthoring}. Do not edit
generated auxiliary files or the PDF directly.

\section{Mandatory Consistency Rules}
\index{AI agent!consistency rules}

\begin{enumerate}
\item Preserve every stable ID. Do not reuse an ID belonging to a retired or
deleted proposal.
\item Use the scope and family codes in
Section~\ref{Section:RequirementIdentifierConvention}. Before using an
undefined scope or family, propose its code, name, boundary, rationale, and
examples; check for semantic overlap and identifier collision; obtain
authorized-human approval; and update the convention, category catalog, this
guide, and the applicable register. Do not rename existing IDs when extending
the convention.
\item Use \textbf{shall} for binding requirements. Use priority terms only in
the Priority field and with the definitions in the glossary.
\item Give every requirement a priority, status, In Version value, use-case
mapping, and test-case mapping.
\item Enter \NoVersion\ when a requirement is not Done and no completed
implementation version applies.
\item Maintain reciprocal mappings among user stories, requirements, use
cases, tests, traceability, verification, and validation.
\item Do not copy one controlled statement into several authoritative
locations. Link to its stable ID.
\item Retain retired requirements and their history. Record the rationale and
replacement when applicable.
\item Use project evidence for factual claims. Label an unsupported suggestion
as a proposal and report the evidence needed to decide it.
\item Record each domain profile as Applicable, Not Applicable, or TBD with an
authorized decision and rationale. Load every Applicable or relevant TBD
profile; do not duplicate project facts in the profile appendix.
\item Before creating or renaming any project-owned file, ask for or confirm
the project \texttt{\textless nameSpace\textgreater}; do not invent it. Prefix
every new project-owned file name with the confirmed namespace.
\item Use TikZ for every diagram. Every TikZ source shall use the directory
\begin{center}
\texttt{TikZ/\textless includingFileBaseName\textgreater/}
\end{center}
and shall be named \path|<nameSpace><DrawingName>.tex|. The directory shall exactly match the base
name of the LaTeX file containing its \verb|\input| command. The including
chapter shall place it in a
numbered figure, add a descriptive caption and stable label, and reference the
figure in its text.
\item Link every document-history entry to the location where the change
occurred.
\item Compile after structural, cross-reference, glossary, index, diagram, or
table changes and visually inspect the affected pages.
\end{enumerate}

\section{Standard AI Workflows}
\label{Section:AIArchitectureWorkflow}
\index{AI agent!architecture workflow}

\subsection{Initialize a Project}
Ask an authorized human for the project file namespace, or confirm an already
configured namespace, before creating or renaming files. Do not invent it.
Replace \verb|\FileNamespace|, project metadata, and bracketed placeholders
only from supplied facts, then apply the namespace to every new project-owned
file.
Identify available domain profiles and ask an authorized human to confirm each
profile's applicability. Record the decisions in
\hyperref[Table:ProjectDomainProfiles]{the project profile table}. Include each
Applicable or TBD profile in the initialization context; remove an
inapplicable profile appendix only after its decision is recorded.
Create a gap list covering purpose, scope, stakeholders, system boundary,
interfaces, constraints, quality attributes, risks, evidence, and approvals.
Resolve consequential gaps before developing detailed artifacts. Work in
dependency order from context and needs through requirements, use cases,
architecture, tests, evidence, and release.

\subsection{Add or Change a Requirement}
Identify the originating need and affected use cases. Draft a singular
verifiable obligation, select its scope, primary family, priority, and status,
and assign a stable ID using
Section~\ref{Section:RequirementIdentifierConvention}. If no defined scope or
family fits, complete the proposal and human-approval procedure before
assigning the ID. Add or update mapped use cases and tests, then update
traceability and verification. Confirm every link in both directions and add a
linked history entry.

\subsection{Add or Change a Use Case}
Identify the actor goal, system boundary, and mapped requirements. Complete one
use-case table with success, alternate, and exception behavior. Add mapped
tests, update requirement mappings and traceability, and confirm that no flow
introduces an unrecorded requirement.

\subsection{Develop or Change Architecture}
Identify the requirements and quality attributes that drive the decision.
Describe alternatives, trade-offs, risks, interfaces, and the selected
structure. Update every affected architectural view and diagram. Confirm that
the proposed design neither omits a driving requirement nor silently adds an
unapproved constraint.

\subsection{Add Tests or Evidence}
Derive each test from mapped requirements and use-case flows. Define repeatable
steps and objective expected results. Link retained evidence only after an
actual execution. Update verification or validation records without converting
an expected outcome into a claimed result.

\subsection{Review and Prepare a Release}
Audit IDs, mappings, statuses, versions, unresolved assumptions, verification,
validation, operational readiness, and release criteria. Report contradictions
and missing evidence. Prepare a release recommendation for human approval; do
not approve the release.

\section{Example AI Requests and Prompts}
\label{Section:AIExamplePrompts}
\index{AI agent!example prompts}
Replace bracketed text with project facts. A prompt does not override this
guide, the controlled project evidence, or human approval boundaries.

\subsection{Initialize the Architecture Book}
\begin{verbatim}
Read tgmAIAgentGuide.tex and use the Initialize a Project workflow
for [Project Name]. Review the supplied project evidence and create
a gap list before editing. Identify missing scope, stakeholders,
interfaces, constraints, quality attributes, risks, evidence, and
approvals. Identify the available domain profiles and ask me to
confirm whether each is Applicable, Not Applicable, or TBD. Read each
Applicable or relevant TBD profile and add its concerns to the gap
list. After the gap list is reviewed, replace only supported
placeholders and develop the document in dependency order. Compile
the complete PDF and report unresolved human decisions.
\end{verbatim}

\subsection{Tailor an Embedded System}
\begin{verbatim}
Read tgmAIAgentGuide.tex and tgmEmbeddedSystemsProfile.tex. Apply the
Embedded Systems Tailoring Profile to [Project Name]. Ask me to
confirm the project namespace, profile applicability, platform class,
hardware revisions, toolchain, interfaces, constraints, applicable
standards, and available evidence. Add embedded concerns to the gap
list before editing. After human review, place project-specific
requirements, architecture, diagrams, tests, and evidence only in
their authoritative core chapters. Do not invent hardware facts or
verification results. Build the complete document, inspect the log,
and visually inspect the affected PDF pages.
\end{verbatim}

\subsection{Create Repository AI Bootstrap Files}
\begin{verbatim}
Read the Repository AI Bootstrap section in tgmAIAgentGuide.tex and
treat that LaTeX chapter as authoritative. Confirm which AI tools
this repository uses and the project file namespace. Create only the
applicable repository discovery files, such as AGENTS.md or
CLAUDE.md. Create SKILL.md only when the repository is intentionally
packaged as, or contains, a reusable skill. Keep every bootstrap file
short: point to tgmManual.tex, tgmAIAgentGuide.tex, the task-routing
table, and tgmDocumentGuide.tex. Do not duplicate architecture rules.
Verify that each file directs architecture-related changes back to
the authoritative LaTeX sources.
\end{verbatim}

\subsection{Develop Requirements}
\begin{verbatim}
Read tgmAIAgentGuide.tex and follow the Requirements context and
workflow. Load every Applicable or TBD domain profile that intersects
the work. From the approved objectives and user stories, propose
singular, feasible, and verifiable requirements. Classify each one,
use shall for the binding obligation, assign one priority and one
status, and map it to at least one use case and test case. Do not
complete requirements work until each applicable profile concern maps
to a new or existing requirement, or to an approved Not Applicable
decision with rationale. Do not mark work Done or assign an
implementation version without evidence.
Update reciprocal mappings, history, and the index; then build and
inspect the affected pages.
\end{verbatim}

\subsection{Develop Use Cases}
\begin{verbatim}
Read tgmAIAgentGuide.tex and follow the Use-Case context and
workflow. Derive use cases from the approved actor goals and mapped
requirements. Give each use case its own table with actors, trigger,
preconditions, main flow, alternate and exception flows, outcome,
postconditions, requirements, and tests. Maintain reciprocal links
and report any flow that requires a new requirement. Build and
inspect the affected pages.
\end{verbatim}

\newpage
\subsection{Develop an Architecture Decision}
\thispagestyle{fancy}
\begin{verbatim}
Read tgmAIAgentGuide.tex and follow the Architecture context and
workflow for [decision or subsystem]. Identify driving requirements,
quality attributes, constraints, interfaces, and risks. Compare
reasonable alternatives and document the trade-off supporting the
proposed choice. Update every affected architecture view and editable
diagram source. Do not present the proposal as approved. Compile the
document and inspect all affected diagrams and pages.
\end{verbatim}

\subsection{Develop Tests and Traceability}
\begin{verbatim}
Read tgmAIAgentGuide.tex and follow the Testing and Evidence context.
For [requirements or use cases], define repeatable tests with inputs,
preconditions, steps, expected results, and required evidence. Update
all reciprocal mappings and the traceability matrix. Do not invent an
execution result or mark verification complete. Compile the document,
check references, and inspect the affected tables.
\end{verbatim}

\subsection{Review Release Readiness}
\begin{verbatim}
Read tgmAIAgentGuide.tex and follow the Release context and workflow
for [release]. Audit requirement status and implementation versions,
traceability, verification, validation, defects, operations, rollback,
known limitations, and approvals. Separate confirmed evidence from
gaps and assumptions. Prepare a release-readiness recommendation for
human review, but do not approve release. Build and inspect the final
document.
\end{verbatim}

\section{Human Decisions and Information Boundaries}
\index{AI agent!human review}
An authorized human shall approve scope, requirements, architecture decisions,
security and privacy posture, safety decisions, evidence sufficiency, waivers,
validation conclusions, and release. The agent shall present alternatives and
evidence without claiming that review has occurred.

Do not provide an AI agent with credentials, secrets, personal data,
production data, export-controlled material, or other information it is not
authorized to process. Before using external sources, confirm that the task
permits network access and that the source may be incorporated. Record material
prompts, assumptions, corrections, and unresolved issues in
\ChapterRef{AIPromptLog}.

\section{Completion Checklist}
\index{AI agent!completion checklist}
\small
\begin{tabularx}{\linewidth}{L{3.5cm}Y}
\caption{AI-Assisted Change Completion Checklist\label{Table:AICompletionChecklist}}\\
\toprule
\textbf{Area} & \textbf{Required Completion Evidence}\\
\midrule
\endfirsthead
\multicolumn{2}{l}{\small\itshape AI-Assisted Change Completion Checklist (continued)}\\
\toprule
\textbf{Area} & \textbf{Required Completion Evidence}\\
\midrule
\endhead
\bottomrule
\endlastfoot
Scope & The requested outcome is complete; unrelated chapters and user changes remain untouched.\\
Domain profiles & Every available profile has an authorized applicability
decision and rationale; every Applicable or relevant TBD profile was included
in the task context, and project facts remain in core chapters.\\
Evidence & New factual statements are supported by supplied or controlled evidence; proposals, assumptions, and unknowns are identified.\\
Language & Binding requirements use shall; priority and status use the controlled terms and definitions.\\
Identifiers & Stable IDs follow the documented scope and family convention, remain unique, and are neither renumbered nor reused; every new scope or family has its definition and authorized-human approval recorded before use.\\
File names & The project namespace was supplied or explicitly confirmed; every newly created project-owned file begins with it, and no namespace was invented.\\
AI bootstrap & Every applicable repository AI discovery file is a thin pointer
to the authoritative LaTeX architecture; tool-mandated filenames are documented
exceptions, and no architecture rules are duplicated in Markdown.\\
Traceability & Every changed requirement, use case, test, verification record, and validation record has complete reciprocal mappings.\\
History & The newest approved document change appears at the top of document history and links to the changed location.\\
Glossary and index & New controlled terms and important concepts are defined and indexed where applicable.\\
Diagram source & Every diagram uses TikZ; each TikZ directory exactly matches its including LaTeX file's base name, each drawing file begins with the project namespace, and the drawing is included, captioned, labeled, and referenced.\\
Build & The complete platform build succeeds; citations, glossary, index, diagrams, and cross-references are regenerated.\\
Log review & No fatal error, undefined reference, duplicate label, missing file, or unresolved layout warning remains.\\
Visual review & Affected PDF pages show readable tables, correct page breaks, current diagrams, working headings, and no clipped or overlapping content.\\
Human review & Decisions requiring authority are clearly identified and are not represented as approved without evidence.\\
\end{tabularx}
\normalsize
